\documentclass[11pt,a4paper]{article}
\usepackage[margin=2.3cm, headsep=14pt]{geometry}
\usepackage{fontspec}
\usepackage{xcolor}
\usepackage{booktabs}
\usepackage{longtable}
\usepackage{array}
\usepackage{fancyhdr}
\usepackage{titlesec}
\usepackage{caption}
\usepackage{enumitem}
\usepackage{microtype}
\usepackage{amsmath}

\renewcommand{\tablename}{Tabla}
\renewcommand{\figurename}{Figura}
\renewcommand{\contentsname}{Contenido}

\definecolor{morado}{HTML}{5E2B8C}
\definecolor{acento}{HTML}{7E57C2}
\definecolor{seccion}{HTML}{EDE7F6}
\definecolor{gristexto}{HTML}{3B3648}

\titleformat{\section}{\Large\bfseries\color{morado}}{\thesection}{0.6em}{}[\vspace{2pt}{\color{acento}\titlerule}]
\titleformat{\subsection}{\large\bfseries\color{acento}}{\thesubsection}{0.5em}{}
\captionsetup{font=small, labelfont={bf,color=morado}, textfont={color=gristexto}}
\setlength{\parindent}{0pt}
\setlength{\parskip}{5pt}
\pagestyle{fancy}
\fancyhf{}
\fancyhead[L]{\small\color{morado}\textbf{Informe de propiedades termof\'isicas}}
\fancyhead[R]{\small\color{gristexto}Intercambiador de Tubos y Coraza}
\fancyfoot[C]{\small\color{gristexto}\thepage}
\renewcommand{\headrulewidth}{0.4pt}

\begin{document}

\begin{titlepage}
\centering
\vspace*{2.2cm}
{\color{acento}\rule{\linewidth}{2pt}}\\[1.2cm]
{\Huge\bfseries\color{morado} Informe de procedencia\\[0.4em]
 y validaci\'on de propiedades}\\[0.8cm]
{\Large\color{acento} Base de datos termof\'isica del software de
 intercambiadores de calor de tubos y coraza}\\[0.5cm]
{\large\color{gristexto} M\'etodo de Kern $\cdot$ Est\'andar TEMA
 Clase C}\\[1.2cm]
{\color{acento}\rule{\linewidth}{1pt}}\\[1.5cm]
\begin{tabular}{rl}
\textbf{Documento generado:} & 17/10/2026 \\
\textbf{Fuente de los datos:} & \texttt{src/core/thermo\_data.py} \\
\textbf{Compilaci\'on:} & LaTeX (XeTeX) con Tectonic \\
\textbf{Cobertura:} & 10 fluidos, 5 materiales, 4 propiedades cada fluido \\
\end{tabular}
\vfill
{\small\color{gristexto} Todas las correlaciones, coeficientes,
 ventanas de validez y fuentes bibliogr\'aficas de la base de datos,
 declaradas de forma completa y trazable.}
\end{titlepage}

\tableofcontents
\newpage

\section{Objeto y alcance}
Este informe declara \textbf{todos los datos} que usa el motor de calculo del
software: de donde proviene cada propiedad termofisica, con que forma
funcional se representa, cuales son sus coeficientes numericos, en que
ventana de temperatura son validas y con que precision reproducen los
valores tabulados de referencia verificados. Los coeficientes que se listan
se leen directamente del codigo fuente (\texttt{thermo\_data.py}) al generar
este documento, de modo que el informe no puede quedar desactualizado
respecto del programa.

Todas las magnitudes estan en unidades SI: $c_p$ en J/(kg\,K), $\mu$ en
Pa\,s, $k$ en W/(m\,K), $\rho$ en kg/m$^3$ y las temperaturas en kelvin (el
programa presenta resultados en grados Celsius). El numero de Prandtl se
calcula como $\mathrm{Pr} = c_p \mu / k$.

El motor de calculo que consume estas propiedades implementa el metodo de
Kern \cite{kern} con la correlacion de conteo de tubos de Serth
\cite{serth}, y el diseno mecanico sigue el estandar TEMA Clase C
\cite{tema} y el ASME BPVC VIII-1 \cite{asme}; la interfaz y las ecuaciones
de dimensionamiento mecanico (19--32) se documentan en Tan y Fok (2006)
\cite{tan-fok}.

\section{Marco metodologico}
\subsection{Formas funcionales}
\textbf{Polinomio centrado} (para $\rho$, $c_p$ y $k$ de liquidos y del
vapor):
\[
  y(T) = c_0 + c_1 (T - T_0) + c_2 (T - T_0)^2 + c_3 (T - T_0)^3
\]
Los polinomios se centran en $T_0$ (temperatura media del rango de ajuste),
lo que reduce la correlacion entre coeficientes y mejora la precision
numerica del ajuste por minimos cuadrados.

\textbf{Viscosidad tipo DIPPR-101} (alcoholes, hidrocarburos, glicol y
amoniaco) \cite{dippr, perry}:
\[
  \mu(T) = \exp\!\left(A + \frac{B}{T} + C \ln T + D\,T\right)
\]

\textbf{Gases} (aire y vapor de agua): densidad por gas ideal
$\rho = P M / (R_u T)$ con $R_u = 8314{,}46$ J/(kmol\,K); para el aire,
$\mu$ y $k$ por la correlacion de Sutherland \cite{sutherland}:
\[
  \mu_{\mathrm{aire}}(T) = 1{,}458\times10^{-6}\,
  \frac{T^{3/2}}{T + 110{,}4}, \qquad
  k_{\mathrm{aire}}(T) = 2{,}646\times10^{-3}\,
  \frac{T^{3/2}}{T + 245{,}4\cdot 10^{-12/T}}
\]

\subsection{Politica de validez}
Cada fluido declara su ventana $[T_{\min}, T_{\max}]$; fuera de ella el
programa lanza un error y \textbf{no extrapola}: la extrapolacion de una
correlacion empirica no es fisica. Ademas, cada propiedad calculada se
valida contra signos fisicos y un rango creible de Prandtl. Las
correlaciones se evaluan a la temperatura media de cada lado, con la
correccion por viscosidad de pared $(\mu/\mu_w)^{0{,}14}$ tipica de
Sieder--Tate.

\section{Procedencia de los datos por fluido}

\begin{longtable}{p{2.6cm}l p{1.5cm} p{6.2cm} p{3.4cm}}
\caption{Procedencia de los datos por fluido}\label{tab:fuentes}\\
\toprule
\textbf{Fluido} & \textbf{Fase} & \textbf{Ventana ($^\circ$C)} &
\textbf{Procedencia del ajuste} & \textbf{Referencias} \\
\midrule
\endfirsthead
\toprule
\textbf{Fluido} & \textbf{Fase} & \textbf{Ventana ($^\circ$C)} &
\textbf{Procedencia del ajuste} & \textbf{Referencias} \\
\midrule
\endhead
Agua & liquido & 0 a 100 & Ajuste por minimos cuadrados a las tablas de liquido saturado/comprimido de IAPWS-IF97 (propiedades de transporte según las formulaciones IAPWS 2008/2011). & \cite{iapws-if97}, \cite{iapws-transporte} \\
Vapor de agua & gas & 100 a 600 & Ajuste a la Tabla A.4 de Incropera (vapor de agua a 1 atm, 380-850 K segun la fuente); densidad por gas ideal. & \cite{incropera} \\
Aire & gas & -73 a 727 & Ajuste a la Tabla A.4 de Incropera (aire seco a 1 atm, 100-3000 K en la tabla); mu y k por la correlacion de Sutherland. & \cite{incropera}, \cite{sutherland} \\
Etanol & liquido & 0 a 75 & Correlaciones tipo DIPPR-101 ajustadas a los datos de Perry's (8.ª ed.) para liquido. & \cite{perry}, \cite{dippr} \\
Metanol & liquido & 0 a 60 & Correlaciones tipo DIPPR-101 ajustadas a los datos de Perry's (8.ª ed.) para liquido. & \cite{perry}, \cite{dippr} \\
Etilenglicol & liquido & 0 a 197 & Correlaciones tipo DIPPR-101 ajustadas a los datos de Perry's / DIPPR (260-470 K) para liquido puro. & \cite{perry}, \cite{dippr} \\
Benceno & liquido & 10 a 80 & Correlaciones tipo DIPPR-101 ajustadas a los datos de Perry's (8.ª ed.) para liquido. & \cite{perry}, \cite{dippr} \\
Tolueno & liquido & -20 a 100 & Correlaciones tipo DIPPR-101 ajustadas a los datos de Perry's (8.ª ed.) para liquido. & \cite{perry}, \cite{dippr} \\
n-Hexano & liquido & -20 a 60 & Correlaciones tipo DIPPR-101 ajustadas a los datos de Perry's (8.ª ed.) para liquido. & \cite{perry}, \cite{dippr} \\
Amoniaco (liquido) & liquido & -40 a 50 & Datos de amoniaco LIQUIDO SATURADO de las tablas ASHRAE / Perry's. & \cite{ashrae}, \cite{perry} \\
\bottomrule
\end{longtable}

\section{Correlaciones y coeficientes por fluido}
Coeficientes tal como estan escritos en \texttt{thermo\_data.py}. $T_0$ es la
temperatura de centrado del polinomio, en K.

\subsection*{Agua}
{\small\color{gristexto} Fase: liquido $\cdot$ Ventana:
 0.0 a 100.0 $^\circ$C (273.15--373.15 K)
 $\cdot$ Fuentes: \cite{iapws-if97}, \cite{iapws-transporte}.}\\[2pt]
\begin{tabular}{llp{7.6cm}}
\toprule
\textbf{Propiedad} & \textbf{Forma funcional} & \textbf{Coeficientes} \\
\midrule
$rho(T)$ & polinomio centrado en $T_0$ = 373.15 K & 958.591, -0.728282, -0.00249653, 5.93186e-06 \\
$cp(T)$ & polinomio centrado en $T_0$ = 373.15 K & 4214.22, 1.32698, 0.0129065, 1.63555e-06 \\
$k(T)$ & polinomio centrado en $T_0$ = 373.15 K & 0.680653, 0.000359543, -7.09323e-06, 1.40325e-08 \\
$\mu(T)$ & DIPPR-101 & -277.891, 10474.5, 44.1717, -0.0534348 \\
\bottomrule
\end{tabular}\\[3pt]
{\footnotesize\itshape Nota de la base de datos: Liquido a 1 atm: hierve exactamente a 100 C. Este software no modela presion de operacion ni cambio de fase, por lo que la ventana de agua liquida TERMINA en 100 C (elegido por el usuario; el ajuste IAPWS subyacente era valido hasta 180 C pero exigiria presurizar el equipo).}
\vspace{0.7em}


\subsection*{Vapor de agua}
{\small\color{gristexto} Fase: gas $\cdot$ Ventana:
 100.0 a 600.0 $^\circ$C (373.15--873.15 K)
 $\cdot$ Fuentes: \cite{incropera}.}\\[2pt]
\begin{tabular}{llp{7.6cm}}
\toprule
\textbf{Propiedad} & \textbf{Forma funcional} & \textbf{Coeficientes} \\
\midrule
$\rho(T)$ & $\rho = P M/(R_u T)$  [gas ideal] & -- \\
$cp(T)$ & polinomio centrado en $T_0$ = 600.00 K & 2017.63, 0.584969, 0.00183989, -6.33368e-06 \\
$k(T)$ & polinomio centrado en $T_0$ = 600.00 K & 0.0421098, 8.29281e-05, 1.38151e-08 \\
$\mu(T)$ & Polinomio & 2.06531e-05, 3.60776e-08, 1.16469e-14 \\
\bottomrule
\end{tabular}\\[3pt]
{\footnotesize\itshape Nota de la base de datos: Vapor sobrecalentado a baja presion. rho por gas ideal (R = 461.5 J/kg.K); por debajo de 373 K a 1 atm el agua es liquida.}
\vspace{0.7em}


\subsection*{Aire}
{\small\color{gristexto} Fase: gas $\cdot$ Ventana:
 -73.1 a 726.9 $^\circ$C (200.00--1000.00 K)
 $\cdot$ Fuentes: \cite{incropera}, \cite{sutherland}.}\\[2pt]
\begin{tabular}{llp{7.6cm}}
\toprule
\textbf{Propiedad} & \textbf{Forma funcional} & \textbf{Coeficientes} \\
\midrule
$\rho(T)$ & $\rho = P M/(R_u T)$  [gas ideal] & -- \\
$cp(T)$ & polinomio centrado en $T_0$ = 600.00 K & 1050.91, 0.225723, 0.000143602, -3.72893e-07 \\
$k(T)$ & Sutherland (ver \S2) & -- \\
$\mu(T)$ & Sutherland (ver \S2) & -- \\
\bottomrule
\end{tabular}\\[3pt]
{\footnotesize\itshape Nota de la base de datos: Aire seco. rho por gas ideal (R = 287.05 J/kg.K); mu y k por Sutherland.}
\vspace{0.7em}


\subsection*{Etanol}
{\small\color{gristexto} Fase: liquido $\cdot$ Ventana:
 0.0 a 75.0 $^\circ$C (273.15--348.15 K)
 $\cdot$ Fuentes: \cite{perry}, \cite{dippr}.}\\[2pt]
\begin{tabular}{llp{7.6cm}}
\toprule
\textbf{Propiedad} & \textbf{Forma funcional} & \textbf{Coeficientes} \\
\midrule
$rho(T)$ & polinomio centrado en $T_0$ = 313.15 K & 772.228, -0.88135, -0.000855147 \\
$cp(T)$ & polinomio centrado en $T_0$ = 313.15 K & 2569.67, 9.23216, 0.0263162 \\
$k(T)$ & polinomio centrado en $T_0$ = 313.15 K & 0.165548, -0.000266076 \\
$\mu(T)$ & DIPPR-101 & 8.30733, 745.101, -3.09521 \\
\bottomrule
\end{tabular}\\[3pt]
{\footnotesize\itshape Nota de la base de datos: Liquido. Punto de ebullicion normal: 78.4 C (351.5 K).}
\vspace{0.7em}


\subsection*{Metanol}
{\small\color{gristexto} Fase: liquido $\cdot$ Ventana:
 0.0 a 60.0 $^\circ$C (273.15--333.15 K)
 $\cdot$ Fuentes: \cite{perry}, \cite{dippr}.}\\[2pt]
\begin{tabular}{llp{7.6cm}}
\toprule
\textbf{Propiedad} & \textbf{Forma funcional} & \textbf{Coeficientes} \\
\midrule
$rho(T)$ & polinomio centrado en $T_0$ = 303.15 K & 781.757, -0.971981, -0.00103094 \\
$cp(T)$ & polinomio centrado en $T_0$ = 303.15 K & 2578.21, 9.64489, 0.0862219 \\
$k(T)$ & polinomio centrado en $T_0$ = 303.15 K & 0.200406, -0.000330078 \\
$\mu(T)$ & DIPPR-101 & -26.1144, 1916.58, 2.13669 \\
\bottomrule
\end{tabular}\\[3pt]
{\footnotesize\itshape Nota de la base de datos: Liquido. Punto de ebullicion normal: 64.7 C (337.8 K).}
\vspace{0.7em}


\subsection*{Etilenglicol}
{\small\color{gristexto} Fase: liquido $\cdot$ Ventana:
 0.0 a 197.0 $^\circ$C (273.15--470.15 K)
 $\cdot$ Fuentes: \cite{perry}, \cite{dippr}.}\\[2pt]
\begin{tabular}{llp{7.6cm}}
\toprule
\textbf{Propiedad} & \textbf{Forma funcional} & \textbf{Coeficientes} \\
\midrule
$rho(T)$ & polinomio centrado en $T_0$ = 323.15 K & 1096.52, -0.682146 \\
$cp(T)$ & polinomio centrado en $T_0$ = 323.15 K & 2517.67, 4.48043 \\
$k(T)$ & polinomio centrado en $T_0$ = 323.15 K & 0.258027, 0.000207156, -2.46019e-06 \\
$\mu(T)$ & DIPPR-101 & -1041.68, 34874.2, 173.839, -0.234198 \\
\bottomrule
\end{tabular}\\[3pt]
{\footnotesize\itshape Nota de la base de datos: Liquido puro (100 \%). Hierve a 197.3 C (470.4 K) a 1 atm; se usa como fluido caliente hasta \textasciitilde{}190 C. Muy viscoso en frio: vigilar Re y dP.}
\vspace{0.7em}


\subsection*{Benceno}
{\small\color{gristexto} Fase: liquido $\cdot$ Ventana:
 10.0 a 80.0 $^\circ$C (283.15--353.15 K)
 $\cdot$ Fuentes: \cite{perry}, \cite{dippr}.}\\[2pt]
\begin{tabular}{llp{7.6cm}}
\toprule
\textbf{Propiedad} & \textbf{Forma funcional} & \textbf{Coeficientes} \\
\midrule
$rho(T)$ & polinomio centrado en $T_0$ = 318.15 K & 851.959, -1.07333 \\
$cp(T)$ & polinomio centrado en $T_0$ = 318.15 K & 1825.45, 4.31919, 0.0111555 \\
$k(T)$ & polinomio centrado en $T_0$ = 318.15 K & 0.137785, -0.000316695 \\
$\mu(T)$ & DIPPR-101 & -20.3514, 1657.46, 1.29529 \\
\bottomrule
\end{tabular}\\[3pt]
{\footnotesize\itshape Nota de la base de datos: Liquido. Congela a 5.5 C y hierve a 80.1 C (353.2 K) a 1 atm.}
\vspace{0.7em}


\subsection*{Tolueno}
{\small\color{gristexto} Fase: liquido $\cdot$ Ventana:
 -20.0 a 100.0 $^\circ$C (253.15--373.15 K)
 $\cdot$ Fuentes: \cite{perry}, \cite{dippr}.}\\[2pt]
\begin{tabular}{llp{7.6cm}}
\toprule
\textbf{Propiedad} & \textbf{Forma funcional} & \textbf{Coeficientes} \\
\midrule
$rho(T)$ & polinomio centrado en $T_0$ = 313.15 K & 847.788, -0.946318 \\
$cp(T)$ & polinomio centrado en $T_0$ = 313.15 K & 1757.03, 3.34971 \\
$k(T)$ & polinomio centrado en $T_0$ = 313.15 K & 0.12612, -0.000322731 \\
$\mu(T)$ & DIPPR-101 & -18.3558, 1393.42, 1.08558 \\
\bottomrule
\end{tabular}\\[3pt]
{\footnotesize\itshape Nota de la base de datos: Liquido. Punto de ebullicion normal: 110.6 C (383.8 K).}
\vspace{0.7em}


\subsection*{n-Hexano}
{\small\color{gristexto} Fase: liquido $\cdot$ Ventana:
 -20.0 a 60.0 $^\circ$C (253.15--333.15 K)
 $\cdot$ Fuentes: \cite{perry}, \cite{dippr}.}\\[2pt]
\begin{tabular}{llp{7.6cm}}
\toprule
\textbf{Propiedad} & \textbf{Forma funcional} & \textbf{Coeficientes} \\
\midrule
$rho(T)$ & polinomio centrado en $T_0$ = 303.15 K & 650.051, -0.892642 \\
$cp(T)$ & polinomio centrado en $T_0$ = 303.15 K & 2288.84, 3.80021 \\
$k(T)$ & polinomio centrado en $T_0$ = 303.15 K & 0.1186, -0.00028 \\
$\mu(T)$ & DIPPR-101 & 15.8121, -349.414, -3.99588 \\
\bottomrule
\end{tabular}\\[3pt]
{\footnotesize\itshape Nota de la base de datos: Liquido. Punto de ebullicion normal: 68.7 C (341.9 K).}
\vspace{0.7em}


\subsection*{Amoniaco (liquido)}
{\small\color{gristexto} Fase: liquido $\cdot$ Ventana:
 -40.0 a 50.0 $^\circ$C (233.15--323.15 K)
 $\cdot$ Fuentes: \cite{ashrae}, \cite{perry}.}\\[2pt]
\begin{tabular}{llp{7.6cm}}
\toprule
\textbf{Propiedad} & \textbf{Forma funcional} & \textbf{Coeficientes} \\
\midrule
$rho(T)$ & polinomio centrado en $T_0$ = 280.00 K & 629.547, -1.42681, -0.0026993 \\
$cp(T)$ & polinomio centrado en $T_0$ = 280.00 K & 4658.1, 6.80046, 0.0222242 \\
$k(T)$ & Sutherland (ver \S2) & -- \\
$\mu(T)$ & DIPPR-101 & -34.8104, 1833.62, 3.46104 \\
\bottomrule
\end{tabular}\\[3pt]
{\footnotesize\itshape Nota de la base de datos: Liquido SATURADO (refrigerante): la presion no es 1 atm sino la de saturacion (p.ej. \textasciitilde{}8.6 bar a 20 C). Hierve a -33.3 C a 1 atm.}
\vspace{0.7em}

\section{Materiales solidos de construccion}
Propiedades de los metales de tubos y coraza \cite{tema}. $k_m$ y $\rho_m$ a
temperatura ambiente; $\varepsilon$ es la rugosidad absoluta y $R_f$ el
factor de ensuciamiento de diseno del estandar TEMA Clase C (1
h\,ft$^2$\,\textdegree F/Btu = 0{,}17611 m$^2$\,K/W; los valores tabulados
corresponden a 0,001 y 0,002 h\,ft$^2$\,\textdegree F/Btu de la Tabla
RGP-2.4).

\begin{tabular}{lrrrr}
\toprule
\textbf{Material} & \textbf{$k_m$ (W/m\,K)} & \textbf{$\rho_m$
 (kg/m$^3$)} & \textbf{$\varepsilon$ (m)} & \textbf{$R_f$
 (m$^2$\,K/W)} \\
\midrule
Cobre & 386.0 & 8933 & 1.5e-06 & 1.761e-04 \\
Acero al carbono & 60.5 & 7854 & 4.6e-05 & 3.522e-04 \\
Acero inoxidable AISI 304 & 14.9 & 7900 & 1.5e-06 & 1.761e-04 \\
Laton Admiralty (70Cu-29Zn-1Sn) & 111.0 & 8530 & 1.5e-06 & 1.761e-04 \\
Aluminio & 237.0 & 2702 & 1.5e-06 & 1.761e-04 \\
\bottomrule
\end{tabular}

\section{Validacion numerica}
Cada punto compara el valor que devuelve la correlacion del programa con el
valor tabulado de referencia verificado (NIST Chemistry WebBook,
CoolProp/REFPROP y correlaciones de referencia publicadas). La desviacion es
$(\text{calculado} - \text{referencia})/\text{referencia}$.

{\small\setlength{\tabcolsep}{4pt}%
\begin{longtable}{@{}l c c r r r l c@{}}
\caption{Validacion frente a valores tabulados de referencia}\label{tab:val}\\
\toprule
\textbf{Fluido} & \textbf{$T$ ($^\circ$C)} & \textbf{Prop.} &
\textbf{Referencia} & \textbf{Programa} & \textbf{Desv.} &
\textbf{Fuente del valor} & \textbf{Conf.} \\
\midrule
\endfirsthead
\toprule
\textbf{Fluido} & \textbf{$T$ ($^\circ$C)} & \textbf{Prop.} &
\textbf{Referencia} & \textbf{Programa} & \textbf{Desv.} &
\textbf{Fuente del valor} & \textbf{Conf.} \\
\midrule
\endhead
Agua & 20 & rho & 998.2 & 997.839 & -0.04\,\% & NIST (IAPWS-95) & alta \\
Agua & 20 & cp & 4184 & 4189.82 & +0.14\,\% & NIST (IAPWS-95) & alta \\
Agua & 20 & k & 0.598 & 0.599308 & +0.22\,\% & NIST / IAPWS & alta \\
Agua & 20 & mu & 0.001002 & 0.00100895 & +0.69\,\% & NIST (Huber 2009) & alta \\
Agua & 60 & rho & 983.2 & 983.348 & +0.02\,\% & NIST (IAPWS-95) & alta \\
Agua & 60 & cp & 4185 & 4181.69 & -0.08\,\% & NIST (IAPWS-95) & alta \\
Agua & 60 & k & 0.651 & 0.654024 & +0.46\,\% & NIST / IAPWS & alta \\
Agua & 60 & mu & 4.660e-04 & 4.636e-04 & -0.50\,\% & NIST (Huber 2009) & alta \\
Agua & 100 & rho & 958.4 & 958.591 & +0.02\,\% & NIST (liq. saturado) & alta \\
Agua & 100 & cp & 4216 & 4214.22 & -0.04\,\% & NIST (liq. saturado) & alta \\
Agua & 100 & k & 0.677 & 0.680653 & +0.54\,\% & NIST / IAPWS & alta \\
Agua & 100 & mu & 2.817e-04 & 2.814e-04 & -0.09\,\% & NIST (Huber 2009) & alta \\
Etanol & 25 & rho & 785.1 & 785.255 & +0.02\,\% & CoolProp 8 & media \\
Etanol & 25 & cp & 2435 & 2437.11 & +0.09\,\% & CoolProp 8 & media \\
Etanol & 25 & k & 0.1635 & 0.169539 & +3.69\,\% & CoolProp 8 & media \\
Etanol & 25 & mu & 0.00108 & 0.00108213 & +0.20\,\% & CoolProp 8 & media \\
Metanol & 25 & rho & 786.3 & 786.591 & +0.04\,\% & NIST (IUPAC) & alta \\
Metanol & 25 & cp & 2535 & 2532.14 & -0.11\,\% & NIST (IUPAC) & alta \\
Metanol & 25 & k & 0.2 & 0.202057 & +1.03\,\% & NIST (IUPAC) & alta \\
Metanol & 25 & mu & 5.440e-04 & 5.464e-04 & +0.44\,\% & NIST (IUPAC) & alta \\
Etilenglicol & 25 & rho & 1110 & 1113.57 & +0.32\,\% & NIST (Zhou-Lemmon) & media \\
Etilenglicol & 25 & cp & 2412 & 2405.66 & -0.26\,\% & REFPROP / DIPPR & media \\
Etilenglicol & 25 & k & 0.252 & 0.251311 & -0.27\,\% & Mebbelli 2021 & media \\
Etilenglicol & 25 & mu & 0.0168 & 0.0169871 & +1.11\,\% & Mebbelli 2021 & media \\
Benceno & 25 & rho & 873.7 & 873.426 & -0.03\,\% & NIST (Thol/Lemmon) & alta \\
Benceno & 25 & cp & 1734 & 1743.53 & +0.55\,\% & NIST (Thol/Lemmon) & alta \\
Benceno & 25 & k & 0.1413 & 0.144119 & +2.00\,\% & NIST (Assael) & alta \\
Benceno & 25 & mu & 6.030e-04 & 6.038e-04 & +0.14\,\% & NIST (Assael) & alta \\
Tolueno & 25 & rho & 862.2 & 861.983 & -0.03\,\% & NIST (Lemmon-Span) & alta \\
Tolueno & 25 & cp & 1701 & 1706.79 & +0.34\,\% & NIST (Lemmon-Span) & alta \\
Tolueno & 25 & k & 0.1304 & 0.130961 & +0.43\,\% & NIST (Assael) & alta \\
Tolueno & 25 & mu & 5.520e-04 & 5.547e-04 & +0.50\,\% & NIST (Assael) & alta \\
n-Hexano & 25 & rho & 654.9 & 654.514 & -0.06\,\% & NIST (Thol/Wang) & alta \\
n-Hexano & 25 & cp & 2273 & 2269.83 & -0.14\,\% & NIST (Thol/Wang) & alta \\
n-Hexano & 25 & k & 0.12 & 0.12 & +0.00\,\% & NIST (Assael) & alta \\
n-Hexano & 25 & mu & 2.980e-04 & 2.955e-04 & -0.83\,\% & NIST (Assael) & alta \\
Amoniaco (liquido) & 25 & rho & 602.9 & 602.761 & -0.02\,\% & NIST (Tillner-Roth) & alta \\
Amoniaco (liquido) & 25 & cp & 4780 & 4788.85 & +0.19\,\% & NIST (Tillner-Roth) & alta \\
Amoniaco (liquido) & 25 & k & 0.47 & 0.4905 & +4.36\,\% & NIST (Tillner-Roth) & alta \\
Amoniaco (liquido) & 25 & mu & 1.270e-04 & 1.309e-04 & +3.10\,\% & NIST (Tillner-Roth) & alta \\
Aire & 25 & rho & 1.184 & 1.1839 & -0.01\,\% & CoolProp (Lemmon 2000) & alta \\
Aire & 25 & cp & 1006 & 1006.12 & +0.01\,\% & CoolProp (Lemmon 2000) & alta \\
Aire & 25 & k & 0.0262 & 0.0261044 & -0.36\,\% & Incropera / CoolProp & alta \\
Aire & 25 & mu & 1.840e-05 & 1.837e-05 & -0.15\,\% & CoolProp (Lemmon 2000) & alta \\
Vapor de agua & 150 & rho & 0.5233 & 0.518835 & -0.85\,\% & NIST (IAPWS-95, real) & alta \\
Vapor de agua & 150 & cp & 1986 & 2006.76 & +1.05\,\% & NIST (IAPWS-95) & alta \\
Vapor de agua & 150 & k & 0.0289 & 0.027876 & -3.54\,\% & NIST / IAPWS & alta \\
Vapor de agua & 150 & mu & 1.420e-05 & 1.427e-05 & +0.52\,\% & NIST (Huber 2009) & alta \\
\bottomrule
\end{longtable}}

\textbf{Lectura de la validacion.} Desviacion maxima en densidad y
capacidad calorifica: 1.05\,\%; en conductividad termica:
4.36\,\%; en viscosidad: 3.10\,\%. La
viscosidad es la propiedad mas sensible a la fuente tabulada; las
discrepancias entre fuentes para un mismo fluido alcanzan el 2--6\,\%
(como se anota en la columna ``Fuente del valor''). Para diseno de detalle
con fluidos de viscosidad critica se recomienda contrastar con NIST
Chemistry WebBook, REFPROP/CoolProp o datos de fabrica.

\section{Conclusiones y limitaciones de uso}
\begin{enumerate}[leftmargin=1.4em]
\item Las correlaciones son validas \textbf{solo} dentro de la ventana
declarada por fluido; fuera de ella el programa rechaza el calculo en lugar
de extrapolar.
\item El agua liquida termina en 100\,\textdegree C (ebullicion a 1 atm). El
ajuste IAPWS subyacente alcanzaba 180\,\textdegree C, pero exigiria
presurizar el equipo y el modelo es monofasico, sin presion de operacion
\cite{iapws-if97}.
\item El amoniaco se modela como liquido \textbf{saturado}: su presion es la
de saturacion (10,03 bar a 25\,\textdegree C), no 1 atm \cite{ashrae}.
\item El aire y el vapor usan gas ideal para $\rho$; frente al valor real del
vapor a 150\,\textdegree C y 1 atm la diferencia es del 0,8\,\%, aceptable
para el alcance docente; para presiones muy alejadas de 1 atm habria que
corregir el factor de compresibilidad \cite{nist}.
\item La viscosidad es la propiedad con mayor dispersion entre fuentes
(2--6\,\% para un mismo fluido) y concentra las mayores desviaciones de la
validacion.
\item Para diseno en regimen laminar, la incertidumbre de los coeficientes de
conveccion domina sobre la de las propiedades: el programa reporta bandas de
hasta $\pm 40\,\%$ sobre $U$ en esos casos.
\end{enumerate}

\section{Bibliografia}
\begin{enumerate}[label={[\arabic*]}, leftmargin=1.6em]
\bibitem{iapws-if97} IAPWS. \emph{R7-97(2012): Revised Release on the IAPWS Industrial Formulation 1997 for the Thermodynamic Properties of Water and Steam}. International Association for the Properties of Water and Steam, Secretariat Erlangen, 1997 (revision 2007; cambios editoriales 2009--2012). Region 1 (agua liquida): $273{,}15\ \mathrm{K} \le T \le 623{,}15\ \mathrm{K}$, $p_{sat}(T) \le p \le 100$ MPa. Articulo de referencia: W. Wagner et al., ``The IAPWS Industrial Formulation 1997 for the Thermodynamic Properties of Water and Steam'', \emph{ASME Journal of Engineering for Gas Turbines and Power}, 122(1):150--184, 2000, DOI 10.1115/1.483186.
\bibitem{iapws-transporte} IAPWS. \emph{Release on the IAPWS Formulation 2008 for the Viscosity of Ordinary Water Substance} y \emph{Release on the IAPWS Formulation 2011 for the Thermal Conductivity of Ordinary Water Substance}. International Association for the Properties of Water and Steam.
\bibitem{perry} D. W. Green (ed.-in-chief) y R. H. Perry (ed. p\^ostumo). \emph{Perry's Chemical Engineers' Handbook}, 8.ª ed. McGraw-Hill, Nueva York, 2008. ISBN 0-07-142294-3; DOI 10.1036/0071422943. Seccion 2 (Physical and Chemical Data, con datos del DIPPR): Tabla 2-313 ``Viscosity of Inorganic and Organic Liquids (Pa\,s)'', con los coeficientes DIPPR por compuesto.
\bibitem{dippr} AIChE Design Institute for Physical Properties (DIPPR). \emph{DIPPR 801 -- Pure Component Data}, ecuacion 101 para la viscosidad de liquido: $\mu(T)=\exp\left(A+B/T+C\ln T+D\,T^{E}\right)$. En este proyecto se usa la forma con $E=1$.
\bibitem{incropera} T. L. Bergman, A. S. Lavine, F. P. Incropera y D. P. DeWitt. \emph{Fundamentals of Heat and Mass Transfer}, 6.ª ed. John Wiley \& Sons, 2007 (4.ª ed. 1996; 7.ª ed. 2011). Apéndice A, Tabla A.4, ``Thermophysical Properties of Gases at Atmospheric Pressure'': aire (100--3000 K) y vapor de agua (380--850 K).
\bibitem{ashrae} ASHRAE. \emph{ASHRAE Handbook -- Fundamentals}, edicion SI 2025. American Society of Heating, Refrigerating and Air-Conditioning Engineers, Atlanta, 2025. Capitulo 30, ``Thermophysical Properties of Refrigerants'': tablas de liquido y vapor saturados del amoniaco (R-717).
\bibitem{sutherland} W. Sutherland. ``The Viscosity of Gases and Molecular Force''. \emph{The London, Edinburgh, and Dublin Philosophical Magazine and Journal of Science}, Serie 5, 36(223):507--531, 1893. DOI 10.1080/14786449308620508.
\bibitem{nist} NIST Chemistry WebBook, NIST Standard Reference Database Number 69. National Institute of Standards and Technology, Gaithersburg, MD. DOI 10.18434/T4D303. Propiedades por ecuaciones de estado de referencia: IAPWS-95 (agua y vapor), Thol/Lemmon/Span (benceno, tolueno, n-hexano), Tillner-Roth (amoniaco), Lemmon \& Jacobsen (aire).
\bibitem{coolprop} I. H. Bell, J. Wronski, S. Quoilin y V. Lemmon. ``Pure and Pseudo-pure Fluid Properties Database for Refrigerants and Fluids'' (CoolProp). \emph{Industrial \& Engineering Chemistry Research}, 53(6):2498--2508, 2014.
\bibitem{eg-referencia} A. Mebbelli, E. Velliadou, M. J. Assael, J. Antoniadis y W. M. Huber. ``Reference Correlations for the Viscosity and Thermal Conductivity of Ethylene Glycol''. \emph{International Journal of Thermophysics}, 42, 2021.
\bibitem{tema} Tubular Exchanger Manufacturers Association, Inc. \emph{Standards of the Tubular Exchanger Manufacturers Association}, Tenth Edition. TEMA, Inc., 2019. Figura N-1.2 (Table of TEMA types) y tablas CB-3.1.3, RCB-4.3, CB-4.4.1, RCB-4.5.2, CB-4.7.1, C-7.1.1, C-11.1 y RCB-4.6.1 de la Seccion 5 (Mechanical Standards TEMA Class R C B).
\bibitem{asme} ASME Boiler \& Pressure Vessel Code, Seccion VIII, Division 1. American Society of Mechanical Engineers. UG-27 (cilindros), UG-32 (cabezas), UG-34 (cubiertas planas).
\bibitem{kern} D. Q. Kern. \emph{Process Heat Transfer}. McGraw-Hill, Nueva York, 1950 (871 pp.). Metodo de dimensionamiento termico del lado coraza con baffles segmentados.
\bibitem{serth} R. W. Serth. \emph{Process Heat Transfer: Principles and Applications}, 1.ª ed. Academic Press (Elsevier), Oxford, 2007. ISBN 978-0-12-373588-1. Correlacion de conteo de tubos con constantes $C_L$ y $C_{TP}$ (ver tambien R. W. Serth y T. Lestina, \emph{Process Heat Transfer: Principles, Applications and Rules of Thumb}, Academic Press, 2014).
\bibitem{tan-fok} F. L. Tan y S. C. Fok. ``An educational computer-aided tool for heat exchanger design''. \emph{Computer Applications in Engineering Education}, 14(2):77--89, 2006. DOI 10.1002/cae.20073. Herramienta educativa de diseno de intercambiadores de tubos y coraza cuya interfaz inspira a este proyecto; las ecuaciones 19--32 de dimensionamiento mecanico citadas por el codigo se refieren a este articulo.
\end{enumerate}

\end{document}
